
\def\PUDcodigo{ECA235}

\def\PUDcodacad{04507.58}

\def\PUDdisciplina{Controle III}

\def\PUDsemestre{Opt}

\def\PUDhoras{80}

%NOVOS ---------------------------------------------------
\def\PUDhorasTeoria{60}
\def\PUDhorasPratica{20}

\def\PUDinterECA{ \hyperref[PUD:ECA220]{Sistemas Lineares (04507.26)}

\hyperref[PUD:ECA233]{Controle I (04507.31)}

\hyperref[PUD:ECA202]{Controle II (04507.40)}

\hyperref[PUD:ECA218]{Identificação de Sistemas (04507.56)}

\hyperref[PUD:EME132]{Inteligência Computacional Aplicada (04507.61)}
}

\def\PUDatuaisECA{---}

\def\PUDrecursos{Projetor multimídia; Quadro branco e pincel.}

%----------------------------------------------------------

\def\PUDcreditos{4}

\def\PUDprerequisito{\hyperref[PUD:ECA202]{Controle II (04507.40)}}

\def\PUDpreacad{\hyperref[PUD:ECA202]{Controle II (04507.40)}}

\def\PUDobjetivos{Aplicar os conceitos de controle;  estudar técnicas de controle avançado.}

\def\PUDementa{Introdução ao controle inteligente. Controle PID. Controle Fuzzy. Controle Ótimo. Controle Preditivo. Controle utilizando Redes Neurais.
%Controle avançado;  controle multivariável;  controle ótimo;  controle preditivo;  controle inteligente.
}

\def\PUDconteudo{ & Introdução ao controle inteligente: controle tradicional, fuzzy, ótimo, preditivo e  neural.
\\
& Controle PID: Fundamentação Matemática. Sintonia de um controlador PID tradicional.
\\
& Controle Fuzzy: Fundamentação matemática. Projeto de Controlador Fuzzy. Controlador híbrido PI-Fuzzy.
\\
& Controle Ótimo: Fundamentação Matemática. LQR e LQI.
\\
& Controle Preditivo: Fundamentação Matemática. GPC.
\\
& Controle utilizando Redes Neurais: Fundamentação Matemática. Projeto de controlador neural.
%&  Modelagem de Sistemas Multivariáveis; 
%\\ 
% &  Elementos Básicos do Controlador Preditivo Baseado em Modelo (CPBM);  Representação de Sistemas Lineares Com Ruído;  Controle Preditivo Generalizado (GPC);  GPC para processos com atrasos, CPBM com Restrições. 
%\\ 
% &  Projeto de Controladores Ótimos; 
%\\ 
% &  Projeto de Controladores Adaptativos. 
}

\def\PUDmetodo{Aulas expositórias que podem ser teóricas e/ou práticas, onde as práticas no laboratório serão marcadas no decorrer da disciplina.}

\def\PUDavalia{A avaliação será feita com aplicação de provas teóricas e/ou práticas, além da possibilidade de inclusão de trabalhos e seminários no decorrer da disciplina.}

\def\PUDbibbasica{ &  \href{www.biblioteca.ifce.edu.br/index.asp?codigo_sophia=63228}{Ogata}, Katsuhiko.  Engenharia de controle moderno.  5a ed.  São Paulo: Pearson.  2011.  824p.  ISBN: 9788576058106. 
\\ 
 &  \href{www.biblioteca.ifce.edu.br/index.asp?codigo_sophia=24226}{Simões}, Marcelo Godoy.  Controle e Modelagem Fuzzy.  Volume único.  Editora Blucher.  São Paulo, 2007.  ISBN: 8521204167. 
\\ 
 &  \href{www.biblioteca.ifce.edu.br/index.asp?codigo_sophia=24230}{Campos}, Mario Cesar M.  Massa.  Controles típicos de equipamentos e processos industriais.  Volume único.  Editora Edgar Blucher.  São Paulo, 2006.  ISBN: 8521203985. }

\def\PUDbibcomplementar{ &  \href{www.biblioteca.ifce.edu.br/index.asp?codigo_sophia=62576}{Maya}, Paulo Álvaro.  Controle essencial.  2º ed.  Pearson Education do Brasil.  2014.  ISBN: 9788543002415.
\\
 & \href{www.biblioteca.ifce.edu.br/index.asp?codigo_sophia=65213}{Aguirre}, Luis Antonio.   Introdução à identificação de sistemas: técnicas lineares e não lineares, teoria e aplicação.  4º ed.   Editora UFMG.  ISBN: 9788542300796.
 %&  Maya, Paulo Álvaro; Leonardi, Fabrizio.  Controle essencial. E-book.  ISBN: 9788576057000. 
\\ 
 & \href{www.biblioteca.ifce.edu.br/index.asp?codigo_sophia=39837}{Lathi}, B. P.  Sinais e sistemas lineares.  2º ed.  Porto Alegre, RS : Bookman, 2007.  856p.  ISBN: 9788560031139.
\\
 & \href{www.biblioteca.ifce.edu.br/index.asp?codigo_sophia=56037}{Shumway-Cook}, Anne; Woollacott, Marjorie H.  Controle Motor: teoria e aplicações práticas.  E-book.  3ª ed.  Manole. 636p.  ISBN: 9788520427477.
\\
 & \href{www.biblioteca.ifce.edu.br/index.asp?codigo_sophia=65382}{Capelli}, Alexandre.  Automação industrial: controle do movimento e processos contínuos.  3º ed.  São Paulo, SP: Érica, 2013.  236p.  ISBN: 9788536501178. }

\def\PUDautor{Samuel Vieira Dias}

\def\PUDdata{2013-05-21}

\def\PUDrevisao{2}

\def\PUDrevisor{Samuel Vieira Dias}

\def\PUDdatarevisao{2019-05-15}

\PUDcorpo
